\documentclass[11pt,reqno]{amsart}
\usepackage[T1]{fontenc}
\usepackage{lmodern}
\usepackage{microtype}
\usepackage{amsmath,amssymb,mathtools}
\usepackage{booktabs}
\usepackage[colorlinks=true,linkcolor=blue,citecolor=blue,urlcolor=blue]{hyperref}
\newtheorem{theorem}{Theorem}[section]
\newtheorem{lemma}[theorem]{Lemma}
\newtheorem{proposition}[theorem]{Proposition}
\newtheorem{corollary}[theorem]{Corollary}
\newtheorem{conjecture}[theorem]{Conjecture}
\theoremstyle{definition}
\newtheorem{definition}[theorem]{Definition}
\newtheorem{remark}[theorem]{Remark}
\newtheorem{computation}[theorem]{Computation}
\newtheorem{citedthm}[theorem]{Cited theorem}
\newtheorem{assessment}[theorem]{Assessment}
\newtheorem{problem}{Problem}
\newcommand{\Z}{\mathbb Z}
\newcommand{\Q}{\mathbb Q}
\newcommand{\N}{\mathbb N}
\newcommand{\Zh}{\widehat{\Z}}
\newcommand{\Zhu}{\widehat{\Z}^{\times}}
\newcommand{\Cl}{\operatorname{Cl}}
\newcommand{\sqf}{\operatorname{sf}}
\newcommand{\ckmin}{ck_{\min}}
\newcommand{\leg}[2]{\left(\frac{#1}{#2}\right)}
\newcommand{\status}[1]{\textup{\textsc{#1}}}
\title[Finite coverings for Erd\H{o}s--Straus and sterile points]{Finite congruence coverings for the Erd\H{o}s--Straus equation and sterile profinite points}
\author{Anonymous}
\date{Draft, 6 October 2026; internal checks only, not externally refereed}
\begin{document}
\begin{abstract}
ABSTRACT-PLACEHOLDER
\end{abstract}
\maketitle

\section{Introduction}\label{sec:intro}

INTRO-PLACEHOLDER

\section{Coverings as clopen classes}\label{sec:frame}

\subsection{The seven families}
We use the families of Elsholtz and Tao \cite[Prop.~1.9]{ET}.  Each is
indexed by a triple $\kappa$ of positive integers subject to a side
condition, and defines a set of residue classes:
\begin{center}
\begin{tabular}{llll}
\toprule
family & parameters & side condition & class\\
\midrule
I1 & $a,d,f$ & $f\mid 4a^2d+1$ & $n\equiv-f \pmod{4ad}$\\
I2 & $a,c,f$ & $(4ac,f)=1$ & $n\equiv-f\ (4ac)$, $n\equiv -c/a\ (f)$\\
I3 & $c,d,f$ & $(4cd,f)=1$ & $n\equiv-f\ (4cd)$, $n^2\equiv-4c^2d\ (f)$\\
I4 & $a,b,e$ & $e\mid a+b$, $(e,4ab)=1$ & $n\equiv-1/e \pmod{4ab}$\\
II1 & $a,b,e$ & $e\mid a+b$, $(e,4ab)=1$ & $n\equiv-e\pmod{4ab}$\\
II2 & $a,d,f$ & $4ad\mid f+1$ & $n\equiv-4a^2d\pmod f$\\
II3 & $a,d,e$ & $(4ad,e)=1$ & $n\equiv-4a^2d-e\pmod{4ade}$\\
\bottomrule
\end{tabular}
\end{center}
On its class each family carries explicit solutions.  With
$\pi^{I}(a,b,c,d)=(abdn,acd,bcd)$ and $\pi^{II}(a,b,c,d)=(abd,acdn,bcdn)$
\cite[\S2]{ET}, the remaining coordinates are
\begin{center}
\begin{tabular}{ll}
\toprule
I1 & $e=(4a^2d+1)/f$, $b=(ne+1)/(4ad)$, $c=(n+f)/(4ad)$\\
I2 & $d=(n+f)/(4ac)$, $b=(na+c)/f$\\
I3 & $a=(n+f)/(4cd)$, $b=\bigl(n+(n^2+4c^2d)/f\bigr)/(4cd)$\\
I4 & $c=(a+b)/e$, $d=(ne+1)/(4ab)$\\
II1 & $c=(a+b)/e$, $d=(n+e)/(4ab)$\\
II2 & $c=(f+1)/(4ad)$, $b=(nc+a)/f$\\
II3 & $c=(n+4a^2d+e)/(4ade)$, $b=ce-a=(n+e)/(4ad)$\\
\bottomrule
\end{tabular}
\end{center}
and the solution is $\pi^I$ for I1--I4 and $\pi^{II}$ for II1--II3.

\begin{lemma}[\status{proved}; {\cite[\S\S2, 10]{ET}}]\label{lem:solve}
Let $\kappa$ be an admissible parameter triple of one of the seven
families and let $n\ge1$ be an integer in its class.  Then the coordinates
above are positive integers, and the resulting $(x,y,z)$ is a solution of
$4/n=1/x+1/y+1/z$ in positive integers.
\end{lemma}

\begin{proof}
The identity $4/n=1/x+1/y+1/z$ follows from the defining relations of
the varieties $\Sigma^I_n$, $\Sigma^{II}_n$ of \cite[(2.1)--(2.9),
(2.13)--(2.21)]{ET}.  Integrality on the class is a direct check; for
example in I3, $n\equiv-f\pmod{4cd}$ gives $a\in\Z$, the second
congruence gives $Y:=(n^2+4c^2d)/f\in\Z$, and $fY\equiv f^2\pmod{4cd}$
with $f$ invertible modulo $4cd$ gives $Y\equiv f\equiv -n$, so $b\in\Z$.
With $\kappa$ fixed, every coordinate is a polynomial in $n$ with
non-negative coefficients that is positive at $n=1$, so all coordinates
are positive for $n\ge1$.
\end{proof}

The positivity statement for all $n\ge 1$ (not only for large $n$) was
checked in the project's review of \cite{PM}; the degrees of
$(x,y,z)$ in $n$ are at most~$4$.

ET call a primitive residue class \emph{solvable by polynomials} if there
are polynomials $P_1,P_2,P_3$ with rational coefficients, taking positive
integer values at all sufficiently large $n$ of the class, with
$4/n=\sum 1/P_i(n)$ identically.  Lemma~\ref{lem:solve} says that every
primitive class contained in one of the seven families is solvable by
polynomials.  We use the converse in the following form.

\begin{citedthm}[{\cite[Prop.~1.9]{ET}}]\label{cit:ET19}
If a primitive residue class is solvable by polynomials, then all
sufficiently large primes in it lie in the union of finitely many classes
of the seven families.
\end{citedthm}

(ET state this separately for Type~I and Type~II solvable classes, with
the families I1--I4 and II1--II3 respectively; every solvable primitive
class is of one of the two types \cite[p.~8]{ET}.)

\subsection{Profinite language}
Let $\Zh=\prod_q\Z_q$ be the profinite completion of $\Z$, with unit group
$\Zhu=\prod_q\Z_q^\times$.  Every class of the table is defined by
congruences modulo an integer $M(\kappa)$ (for I3, modulo $4cdf$); we
write $\Cl(\kappa)\subset\Zh$ for the set of $x\in\Zh$ satisfying them.
It is open and closed (\emph{clopen}), and an integer $n$ lies in
$\Cl(\kappa)$ exactly when it lies in the class.  We call $\Cl(\kappa)$
an \emph{ET class}.

\begin{definition}\label{def:sterile}
A point $x\in\Zhu$ is \emph{sterile} if it lies in no ET class.  For a
subfamily $\mathcal K$ of parameter triples, $x$ is
\emph{$\mathcal K$-sterile} if it lies in no $\Cl(\kappa)$,
$\kappa\in\mathcal K$.
\end{definition}

Let $\mathcal P$ be an infinite set of primes and $\mathcal P'$ its set of
accumulation points in $\Zh$.

\begin{lemma}[\status{proved}]\label{lem:accum}
$\mathcal P'$ is a compact subset of $\Zhu$.  If $\Sigma\subset\Zhu$ is
clopen, defined modulo $M$ (a union of reduced classes modulo $M$), and
$\mathcal P(\Sigma)$ is the set of primes $p\nmid M$ with $p\bmod M$ in
$\Sigma\bmod M$, then $\mathcal P(\Sigma)'=\Sigma$.
\end{lemma}

\begin{proof}
If $x\in\mathcal P'$, then for each prime $q$ and each $k$ infinitely many
$p\in\mathcal P$ satisfy $p\equiv x\pmod{q^k}$; all but one of them are
prime to $q$, so $x_q\in\Z_q^\times$.  Derived sets in the compact metric
space $\Zh$ are closed, hence compact.  For the second statement, a
point of $\mathcal P(\Sigma)'$ is congruent modulo $M$ to some element of
$\mathcal P(\Sigma)$, so it lies in $\Sigma$.  Conversely every
neighbourhood $x+Q\Zh$ ($M\mid Q$) of $x\in\Sigma$ is a reduced class
modulo $Q$ and contains infinitely many primes by Dirichlet's theorem;
they lie in $\mathcal P(\Sigma)$.
\end{proof}

The sets of primes studied below are all of the form $\mathcal
P(\Sigma)$.  The prime $p$ itself is never a unit of $\Z_p$, which is why
we work with accumulation points rather than with the primes as points of
$\Zh$.

\begin{proposition}[\status{proved}]\label{prop:compact}
Let $\mathcal P$ be an infinite set of primes.
\begin{enumerate}
\item[(a)] If no point of $\mathcal P'$ is sterile, then finitely many ET
classes contain all but finitely many $p\in\mathcal P$.  For those $p$,
Lemma~\ref{lem:solve} gives an explicit solution of
$4/p=1/x+1/y+1/z$.
\item[(b)] If some $x\in\mathcal P'$ is sterile, then no finite set of ET
classes contains all but finitely many $p\in\mathcal P$.  Moreover no
finite set of primitive residue classes, each solvable by polynomials,
contains all but finitely many $p\in\mathcal P$.
\end{enumerate}
Hence all but finitely many elements of $\mathcal P$ lie in finitely many
primitive classes solvable by polynomials if and only if $\mathcal P'$
contains no sterile point.
\end{proposition}

\begin{proof}
(a) The ET classes are open and cover the compact set $\mathcal P'$, so
finitely many of them cover it; let $U$ be their union.  If
$\mathcal P\setminus U$ were infinite, it would have an accumulation
point, which lies in $\mathcal P'\subset U$; but $U$ is open and misses
$\mathcal P\setminus U$, a contradiction.

(b) Removing finitely many elements from $\mathcal P$ does not change
$\mathcal P'$.  If finitely many ET classes contain a cofinite subset of
$\mathcal P$, their union is closed and therefore contains $x$, which
contradicts sterility.  Next let $C_1,\dots,C_s$ be primitive classes
solvable by polynomials whose union contains a cofinite subset of
$\mathcal P$.  Viewed in $\Zh$, their union is closed, so $x\in C_j$ for
some $j$.  By Cited Theorem~\ref{cit:ET19} there are ET classes
$K_1,\dots,K_t$ containing all large primes of $C_j$.  Their union is
closed and misses $x$, so some neighbourhood $V=x+Q\Zh\subset C_j$
misses it.  As $x$ is a unit, $V$ is a reduced class modulo $Q$, and by
Dirichlet's theorem it contains infinitely many primes.  These lie in
$C_j$ and in no $K_i$, a contradiction.  The final equivalence combines
(a), (b) and Lemma~\ref{lem:solve}.
\end{proof}

\begin{remark}[Direction of the implications]\label{rem:direction}
The equivalence concerns the \emph{existence} of a sterile point of
$\mathcal P'$.  For a \emph{given} point $x$ only one direction holds:
if $x$ is sterile there is no finite covering, but if $x$ lies in some
ET class another point of $\mathcal P'$ may still be sterile.  An earlier
internal draft of \cite{PM} said ``equivalently'' here; this was false
and was corrected in review.  A weaker statement is available for every
explicit point: if $x\in\mathcal P'$ lies in no ET class of modulus at
most $Y$, then every finite ET covering of a cofinite subset of $\mathcal
P$ uses a class of modulus greater than $Y$.  Indeed the classes of
modulus $\le Y$ in the covering miss a neighbourhood of $x$, which
contains infinitely many elements of $\mathcal P$.
\end{remark}

\begin{remark}[Heights]\label{rem:heights}
The same compactness argument works for any subfamily $\mathcal K$ and
any function $\kappa\mapsto\operatorname{ht}(\kappa)$.  Let
$\mathcal K_X$ be the triples of height at most $X$ and
$U_X=\Sigma\setminus\bigcup_{\kappa\in\mathcal K_X}\Cl(\kappa)$ for a
compact $\Sigma$.  The $U_X$ are closed and decreasing.  If each is
nonempty, the finite intersection property gives a point of $\bigcap_X
U_X$, which is $\mathcal K$-sterile.  So either some finite subfamily of
bounded height covers $\Sigma$, or $\Sigma$ has a $\mathcal K$-sterile
point.  The set of $\mathcal K$-sterile points of $\Sigma$ is
$\bigcap_XU_X$, a closed set.
\end{remark}

\section{Square points and the odd-square obstruction}\label{sec:square}

\begin{proposition}[\status{proved}; known]\label{prop:square}
Let $x\in\Zhu$ be a \emph{square point}: $x_q\in(\Z_q^\times)^2$ for every
prime $q$.  Then $x$ is sterile.
\end{proposition}

\begin{proof}
For $x=1$ this is elementary: if $1\in\Cl(\kappa)$ then the integer $n=1$
lies in the class, and Lemma~\ref{lem:solve} gives positive integers with
$4=1/x+1/y+1/z\le3$.  In general, suppose $x\in\Cl(\kappa)$ with
$\Cl(\kappa)$ defined modulo $M$.  For each $q\mid M$ the component
$x_q$ is the square of a unit, so by the Chinese remainder theorem there
is an integer $m$ prime to $M$ with $m^2\equiv x\pmod M$.  The primitive
class $m^2\bmod M$ is then contained in the class of $\kappa$ and is
solvable by polynomials (Lemma~\ref{lem:solve}).  This contradicts the
theorem of Mordell and Schinzel that a primitive class whose residue is a
perfect square is not solvable by polynomials
\cite{Mordell,Schinzel}, \cite[p.~8]{ET}.
\end{proof}

The last step is the classical obstruction; ET also derive it from their
vanishing result at odd squares \cite[Prop.~1.6]{ET}, which goes back to
Schinzel and Yamamoto \cite{Yamamoto}.  The input is one quadratic
character computation; Bright and Loughran \cite{BL} identify it with the
Brauer--Manin obstruction of the Erd\H{o}s--Straus surfaces.

Call a prime $p$ \emph{Mordell-hard} if $p\bmod 840$ is a square,
$p\bmod 840\in\{1,121,169,289,361,529\}$.  Every primitive class modulo
$840$ that is not a square is solvable by polynomials (Mordell; see
\cite[p.~8]{ET}).  The set $\Sigma_{\rm hard}\subset\Zhu$ of square
classes modulo $840$ contains the point $1$, so by
Propositions~\ref{prop:compact}(b) and~\ref{prop:square} the Mordell-hard
primes admit no finite covering by polynomial identities.  This is the
classical form of the obstruction, and it is the reason why every finite
covering statement below restricts the hard primes further by a
\emph{non-residue} condition at some additional prime $r$; the resulting
target sets contain no square point.  Whether they contain other sterile
points is the question of this note.

\subsection{The procedure-level form}
The project's working note \cite{PS} proves a version of the odd-square
obstruction for procedures.  We only state it informally and refer to
\cite[\S\S1--2]{PS} and \cite{PObs} for the model of computation.
A \emph{witness-producing program} takes a prime $p$, performs ring
operations, floor divisions, eventual-sign comparisons, full
factorisations and loops over the lists so produced, and outputs a
solution or FAIL; it is \emph{correct} if every solution it outputs is
genuine, and \emph{bounded} if the number of steps is bounded
independently of $p$.  For $P\in\Z[X]$ a point $q^*\in\Zh$ is
\emph{square-mimicking} for $P$ if $P(q^*_\ell)$ is the square of a unit
of $\Z_\ell$ for every prime $\ell$.

\begin{citedthm}[Formal odd-square principle {\cite[Theorem~C,
Corollary~C1]{PS}}; \status{conditional} on Hypothesis~H]\label{cit:C}
Let $P=uX+v$ be linear and $q^*$ square-mimicking for $P$ and
nondegenerate in the sense of \cite[\S1.2]{PS}.  If a correct witness-producing program has a defined and
finite formal run at $(P,q^*)$, then it outputs FAIL at every admissible
$q$ (\status{proved}); under Schinzel's Hypothesis~H \cite{SS} for an
explicit finite family there are infinitely many admissible $q$, and
$p=P(q)\equiv1\pmod{24}$.  In particular, under H for one explicit
finite family, every correct bounded witness-producing program fails for
infinitely many primes $p\equiv1\pmod{24}$.
\end{citedthm}

A finite list of ET classes, each tested by a congruence and followed by
the identity of Lemma~\ref{lem:solve}, is such a bounded program, and
its failure at the rational square-mimicking point $1$ is
Proposition~\ref{prop:square} for $x=1$, unconditionally.  Cited Theorem~\ref{cit:C} does
not reprove the unconditional statement for general square points; it
extends the obstruction from congruence coverings to all bounded
procedures, at the price of Hypothesis~H.  \emph{Novelty.} Theorem~C is
the procedure-level, H-conditional generalisation of known obstructions
(Mordell and Schinzel on square classes, ET Prop.~1.6 and the remark on
\cite[p.~6]{ET} that ``a finite set of covering congruence strategies''
must fail, Yamamoto, Bright--Loughran); its unconditional instances are
known in substance.  The project's novelty audit rates it as a partial
novelty only.

\begin{remark}[\status{Assessment}]\label{rem:notsquare}
The target sets of \S\S\ref{sec:typeI}--\ref{sec:r17} require $x_r$ to be a
non-square for some prime $r$, so they contain no point that is
square-mimicking for $P=X$, and neither Proposition~\ref{prop:square} nor
Cited Theorem~\ref{cit:C} applies to them.  The candidate sterile points
below (the sign point $\hat x_9$, the point $x^*$ for $r=13$, the
$17$-generic line) are therefore not explained by the odd-square
obstruction.  If they are sterile, the reason is a different, rigidity
type of obstruction (\S\ref{sec:typeI}); this is a heuristic reading,
not a theorem.
\end{remark}

\section*{Acknowledgements and provenance}
PROVENANCE-PLACEHOLDER

\begin{thebibliography}{99}
\bibitem{BL} M.~Bright and D.~Loughran, Brauer--Manin obstruction for
Erd\H{o}s--Straus surfaces, Bull.\ Lond.\ Math.\ Soc.\ 52 (2020), 746--761;
arXiv:1908.02526v2.
\bibitem{ET} C.~Elsholtz and T.~Tao, Counting the number of solutions to the
Erd\H{o}s--Straus equation on unit fractions, J.\ Aust.\ Math.\ Soc.\ 94
(2013), 50--105; arXiv:1107.1010v6.
\bibitem{Lenstra} H.~W.~Lenstra Jr., Divisors in residue classes, Math.\
Comp.\ 42 (1984), 331--340.
\bibitem{MD} S.~Mihnea and B.~C.~Dumitru, Further verification and empirical
evidence for the Erd\H{o}s--Straus conjecture, arXiv:2509.00128 (2025).
\bibitem{Mordell} L.~J.~Mordell, \emph{Diophantine Equations}, Academic
Press, 1969, Chapter~30 (not accessed; cited via \cite{ET}).
\bibitem{Notes} \emph{Erd\H{o}s--Straus research notebook},
\texttt{notes.md}, this repository, \S\S26, 36, 48; working notes, 2026.
\bibitem{PObs} Anonymous, Under Hypothesis~H, formally generic primes
block the pointwise signed-graph approach to the Erd\H{o}s--Straus
conjecture, research note, this repository
(\texttt{paper/pointwise-obstruction.tex}), 2026.
\bibitem{PS} Companion note in this repository, \texttt{POINTWISE\_SIZE.md}
(Theorems M and C); internally reviewed in
\texttt{reviews/pointwise-size-step1-review.md}.
\bibitem{PT1} Companion note in this repository,
\texttt{POINTWISE\_TYPEI.md}; internally reviewed in
\texttt{reviews/pointwise-typei-review.md}.
\bibitem{PT2} Companion note in this repository,
\texttt{POINTWISE\_TYPEI2.md}; internally reviewed in
\texttt{reviews/pointwise-typei2-review.md}.
\bibitem{PT3} Companion note in this repository,
\texttt{POINTWISE\_TYPEI3.md}; internally reviewed in
\texttt{reviews/pointwise-typei3-review.md}.
\bibitem{PM} Companion note in this repository,
\texttt{POINTWISE\_MORDELL.md}; internally reviewed in
\texttt{reviews/pointwise-mordell-review.md}.
\bibitem{PM17} Companion note in this repository,
\texttt{POINTWISE\_MORDELL17.md}; internally reviewed in
\texttt{reviews/pointwise-mordell17-review.md}.
\bibitem{Salez} S.~E.~Salez, The Erd\H{o}s--Straus conjecture: new modular
equations and checking up to $N=10^{17}$, arXiv:1406.6307 (2014).
\bibitem{Schinzel} A.~Schinzel, On sums of three unit fractions with
polynomial denominators, Funct.\ Approx.\ Comment.\ Math.\ 28 (2000),
187--194 (not accessed; cited via \cite{ET}).
\bibitem{SS} A.~Schinzel and W.~Sierpi\'nski, Sur certaines hypoth\`eses
concernant les nombres premiers, Acta Arith.\ 4 (1958), 185--208.
\bibitem{Yamamoto} K.~Yamamoto, On the Diophantine equation
$4/n=1/x+1/y+1/z$, Mem.\ Fac.\ Sci.\ Kyushu Univ.\ Ser.\ A 19 (1965), 37--47
(not accessed; cited via \cite{ET}).
\end{thebibliography}
\end{document}
